\documentclass[11pt]{article}
\usepackage{lmodern}
\usepackage[T1]{fontenc}
\usepackage[letterpaper,margin=0.8in]{geometry}
\usepackage{xcolor}
\usepackage{booktabs}
\usepackage{tabularx}
\usepackage{enumitem}
\usepackage{array}
\usepackage[hidelinks]{hyperref}

\setlist{itemsep=1pt,topsep=3pt}
\newcommand{\h}[1]{\texttt{#1}}
\newcolumntype{L}{>{\raggedright\arraybackslash}X}
\newcolumntype{P}[1]{>{\raggedright\arraybackslash}p{#1}}

\title{larpix-lib on PACMAN\\[2pt]
  \large LArPix control software:  architecture, setup, bench results and status}
\author{Maintained by the PACMAN agents (augie)}
\date{2026-11-05}

\begin{document}
\maketitle
\sloppy

\section{Purpose and scope}

\h{larpix-lib} (package \h{larpix\_control}) is the new LArPix control software for PACMAN with the zipline firmware
and software (\h{zipline.tex}).  Since 2026-11-05 it is a tool Michael supports officially, no longer an experiment,
and this record moved back into pacman\_docs \h{agents/} beside \h{zipline.tex}.  The goal is for \h{larpix-lib} to
replace (1) the control path, including readback for control, and (2) the data server;  the focus is (1).  The ASICs
themselves (bench setup, registers, TX settings) are recorded in \h{larpix.tex}.

Scope (Michael 2026-11-05):  a single Hydra strand on one UART channel, taken far;  more than one channel is out of
scope for now (the binding of io groups and channels to strands is postponed), but the design must scale to the whole
system.  Every PACMAN carries precisely one ASIC type, so the ASIC version is an attribute of the PACMAN.

\section{Repository and environment}

\begin{itemize}
\item \h{mulhearn/larpix-lib}, cloned in \h{/agents/cowork/larpix-lib} (shared repository;  Michael pushes).  Work is
  on the branch \h{working} (no releases yet;  \h{main} follows \h{origin/main}).  At 3f0bbaa:  PEP 8 names
  (e1e2dbb), David Yang's correlator (fea7566), the closed loop (322c119), board control and
  \h{demos/hello\_asic.py} (2c3e19f), the register dump and 16000-word requests (4f84da5), David Yang's data
  packet layout (da768fc), rails, the data spigot, pedestals and the occupancy view (3f0bbaa).  Needs PACMAN
  software 3.6.3 (requests over 100 words per UART).
\item Python environment:  conda environment \h{larpix} per user, an editable install of the clone;  \h{make test}
  runs the unit tests (46, OK).
\item Names (Michael 2026-11-05):  Python classes in CapWords (PEP 8), functions, methods, modules snake\_case,
  constants UPPER\_CASE.  ``PACMAN'' names only the real board or an accurate emulation;  toys and generic contracts
  say ``board''.
\end{itemize}

\section{Architecture}

\h{larpix\_control} separates what varies independently:  the ASIC version, what each configuration step means, the
network topology, register access to the ASICs, and control of the board they sit on.  Version-specific content
lives in YAML, not code.  The contracts are abstract base classes in \h{common.interfaces}.

\begin{center}
\small
\begin{tabularx}{\linewidth}{@{}P{0.11\linewidth} P{0.24\linewidth} L@{}}
\toprule
Layer & Module & Role \\
\midrule
ASIC spec & \h{asic.\allowbreak asic\_spec} & One ASIC version from YAML (\h{config/asics/larpix\_v3.yaml}):  register
  space, config and data packet layouts.  Field updates to register writes (\h{build\_write\_list}), registers to
  config packets, the read side (\h{valid\_config\_read\_response}, \h{parse\_chip\_address\_value}) and data
  packets (\h{valid\_data\_packet}, \h{parse\_data\_packet\_fields}).  No I/O. \\
Fragments & \h{asic.\allowbreak fragment\_lib} & Named configuration steps in YAML (\h{config/fragments/}) with
  external parameters and per-version overrides. \\
ASIC config & \h{AsicConfigIface}, \h{asic.\allowbreak asic\_config} & The operations a network needs (chip ID, I/O
  init, input and output enables), and any named fragment (\h{apply\_fragment}), one step up from register access.
  \h{AsicConfig} evaluates a fragment, turns the field updates into \h{RegisterOp} writes on its io channel and
  transmits them. \\
Topology & \h{hydra.\allowbreak strand} & One Hydra network on one UART, from YAML (\h{config/hydra/});  configures
  it from the root outward.  Uses only \h{AsicConfigIface}. \\
Register access & \h{AsicRegisterIface} & \h{transmit(ops, fire\_and\_forget)}:  \h{RegisterOp}s
  (\h{OpType} WRITE or READ, \h{nrep} repeats, io channel, chip, register, value) in order;  returns \h{OpReplies}:
  per register (\h{RegisterKey}) the \h{WriteReply} (requested, acknowledged) and \h{ReadReply} (requested,
  acknowledged, majority value, received, discrepant), and an overall \h{complete}. \\
Board control & \h{BoardControlIface} & One readout board (one io group):  \h{ping}, \h{get\_tile},
  \h{reset} (FULL or INTERNAL, by tile) with \h{wait\_on\_reset}, \h{power\_on}/\h{power\_off} (by tile),
  \h{rail\_off}/\h{rail\_on} (VDDA or VDDD, by tile), \h{uart\_enable}/\h{uart\_disable} (by io channel).
  Pokes return at once, without replies. \\
Data & \h{AsicDataIface} & Everything the ASICs on one board send, as it arrives:  \h{start\_data},
  \h{read\_data(timeout\_ms)} (a list of \h{RxWord}:  io channel, 64-bit word, the board's receive time stamp),
  \h{stop\_data};  read regularly (a stalled subscriber makes \h{pacman\_dataserver} quit). \\
PACMAN & \h{pacman.\allowbreak pacman} & \h{Pacman} implements all three:  one PACMAN from one network file entry;
  ASIC words and board register writes through \h{pacman\_cmdserver} (each acknowledged), replies from
  \h{pacman\_dataserver} matched by the reply correlator (\h{asic.\_correlator}).  \h{load\_network} builds one per
  entry;  \h{sync\_pacman\_id} sets the PACMAN's own ID to its io group;  up to 16000 words per request;  the data
  spigot on a SUB socket of its own;  rails by DAC (0 or the \h{pacman\_menu} nominal). \\
Toy & \h{common.\allowbreak board\_toy} & \h{BoardToy} implements all three without hardware:  registers in
  memory, resets, power, rails and UART enables with their effect on the chips;  no data (simulation is for David
  Yang's simulator). \\
\bottomrule
\end{tabularx}
\end{center}

A configuration runs:  \h{Strand.init\_network} $\to$ \h{AsicConfig} operation $\to$ fragment $\to$ field updates
$\to$ register writes $\to$ \h{RegisterOp}s $\to$ \h{Pacman.transmit} $\to$ config packets $\to$ \h{pacman\_cmdserver};
replies return through \h{pacman\_dataserver} to the correlator and the \h{OpReplies}.

\subsection{Design decisions (Michael 2026-11-05)}

\begin{itemize}
\item Chip identity is \h{(io\_chan, chip\_id)}:  chip IDs are unique only within an io channel.  Each board owns a
  configured range of io channels (\h{first\_io\_chan}, \h{uart\_mask});  the upper bits may encode the PACMAN by
  configuration, never checked.  The wire channel is the io channel's low byte (UART + 1);  the UART is
  \h{io\_chan - first\_io\_chan}, so no ``+1'' appears in our code.
\item One board per io group:  the board knows its io group, so board control takes tiles and io channels only.
\item Replies:  one persistent SUB socket per PACMAN, subscribed only during a transmission (an idle subscriber cannot
  stall the data server) and drained at its start (no stale reply counts as fresh).  The race between the
  subscription and the first reply is accepted until the 3.7 XPUB confirmation.
\item The timeout is a property of each implementation (\h{timeout\_ms} in the network file, 200~ms), bounding the
  acknowledgements and the wait for replies.  Missing acknowledgements and replies are reported, not raised.
  \h{fire\_and\_forget} skips collecting replies;  acknowledgements are still counted.
\item Reads repeated (\h{nrep} or listed again) are collated:  the majority value (the last received on a tie) and the
  number discrepant.
\item PACMAN registers, power values and the reset pokes are private to \h{Pacman}, from \h{pacman\_vspace.cc}.  The
  code there sends the internal reset at 0x00100410 and the full reset at 0x00100420;  its header comment has the two
  swapped.
\item Interface first, in Python:  the interface step reviews generated \h{.pyi} stubs (baseline committed on the
  interface branch, the draft in the stubs, propagated into the \h{.py} once accepted);  the stubs never reach
  \h{working}.
\end{itemize}

\paragraph{Later decisions (Michael 2026-11-05, evening).}
\begin{itemize}
\item Interfaces are reviewed in Python as generated \h{.pyi} stubs (\h{pyi-baseline.py}):  a baseline of the
  modules a topic changes committed on the interface branch, the draft in the stubs, then propagated into the
  \h{.py};  configuration schemas (YAML) reviewed with them.
\item The data spigot is its own ABC, not board control;  the ASIC configuration for a task is a fragment applied
  with \h{apply\_fragment}.
\item Data packet fields that diagnostics overlay on the timestamp carry \h{enabled\_by} (the FIFO counts, the
  first-packet marker) or \h{disabled\_by} (the timestamp);  parsing defaults to nothing enabled, and this version
  enables none.
\item Error replies of \h{pacman\_cmdserver}:  one code, 0xEEEE, the cause in its log.
\item The HDF5 data server is a separate component (nothing in \h{larpix\_control} depends on it), tabled;  the
  occupancy monitor belongs to \h{larpix-lib}.
\end{itemize}

\section{Network file}

\h{config/network/*.yaml}, one entry per PACMAN:

\begin{center}
\small
\begin{tabularx}{\linewidth}{@{}lL@{}}
\toprule
Field & Meaning \\
\midrule
\h{io\_group} & The PACMAN;  set as its own ID by \h{sync\_pacman\_id}. \\
\h{hostname}, \h{port} & \h{pacman\_cmdserver} (5555). \\
\h{data\_port} & \h{pacman\_dataserver};  optional, default 5556. \\
\h{timeout\_ms} & Acknowledgement and reply timeout (200). \\
\h{first\_io\_chan}, \h{uart\_mask} & The board's io channels:  \h{first\_io\_chan} + $u$ for each bit $u$;  the
  low byte of each must be its PACMAN channel. \\
\h{uarts\_per\_tile} & 4 on every revision (1v5:  10 tiles, 40 UARTs;  1v4:  8 tiles, 32 UARTs). \\
\bottomrule
\end{tabularx}
\end{center}

The bench file is \h{single\_bench.yaml}:  the 1v5 at 10.42.0.175, io group 1, \h{first\_io\_chan} 0x001,
\h{uart\_mask} 0xF (tile 0, io channels 1--4).

\section{Hello ASIC}

\h{demos/hello\_asic.py}, from the repo top, takes a PACMAN just booted or powered up to a configured, answering strand:
\begin{enumerate}
\item Ping until \h{pacman\_cmdserver} answers (up to \h{--wait}, 120~s), then \h{sync\_pacman\_id}.
\item On the strand's tile:  UART off, \h{power\_on}, 50~ms, full reset, \h{wait\_on\_reset} (1~ms).
\item Configure the strand (\h{Strand} $\to$ \h{AsicConfig} $\to$ \h{Pacman}), recording every register written;
  UART on.
\item Read back every register written and each chip's ID register;  PASS or FAIL, and the exit status.
\end{enumerate}
Options:  \h{--network}, \h{--io\_group}, \h{--io\_chan}, \h{--hydra}, \h{--asic}, \h{--wait}, \h{-v} (the
configuration trace and the board's warnings).  No RX buffer enables are written:  the correlator ignores everything
but config replies.  It replaces the earlier \h{agents/bringup.sh} and \h{pacman\_util.py}.

The settles come from a sweep on the 1v5 (tile 0 off 1~s before each run, 10 runs per step):  a full reset passes
with no settle at all (the reset is about 13~\textmu{}s, shorter than the next round trip), and the power settle
passes from 10~ms and fails at 5~ms and below.  50~ms and 1~ms keep margins of 5$\times$ and about 75$\times$;  the
1v4 and other tiles are not measured.

\section{Pedestals and occupancy}

Pedestals copy larpix-control's \h{pedestal\_qc.py} (larpix-10x10-scripts), in two fragments:
\h{pedestal\_trigger} (periodic and rolling periodic trigger, periodic reset, no hit veto, 100000 trigger cycles,
4096 reset cycles, register 128 complete, all 64 CSAs on) and \h{pedestal\_channels} (trigger and channel masks
cleared);  \h{pedestal\_stop} turns the trigger off again.  The cycle counts are fixed bytes in the fragment
(fragments cannot split a number across a byte array).  \h{demos/pedestal.py} writes both to every chip of the
strand with the tile's VDDA off (as larpix-control), counts the data and stops the trigger (\h{--keep} leaves it).

\h{demos/occupancy.py} is a live strand view of one io channel:  a row per chip (hydra file order) and an ``other''
row, one character per channel (the digit $n$ for at least $10^n$ hits/s, a dot under 1/s, blank for none), per row
hits/s and mean ADC, refreshed each second over a sliding window.  With the rolling trigger every channel sits at
about 1.6 hits/s, so each cell shows 0.

\section{Results}

\begin{center}
\small
\begin{tabularx}{\linewidth}{@{}llL@{}}
\toprule
Date & Board & Result \\
\midrule
2026-11-02 & 1v5 & First configuration by \h{larpix-lib} (\h{demos/init\_network.py} from felix, through
  \h{pacman\_cmdserver}), with the input enables left out:  13 writes.  Then a menu hello without the configuration
  root:  the reply exact, 3 of 3.  Before it, after a full reset, hello had no reply. \\
2026-11-02 & 1v5 & With the input enables written (126 = 0x02, POSI1 only;  14 writes):  the reply exact, 3 of 3;
  read back 124 = 0x00, 125 = 0x01, 126 = 0x02, 248 = 0x07.  UART 0 drives POSI1 (west), as the port map has it. \\
2026-11-02 & 1v5 & The full configuration sequence (fce6523, ``ASIC config working at test bench''):  the reply
  exact, 3 of 3;  io channels 2--4 unanswered, channel 1 exact. \\
2026-11-03 & 1v5 & The whole chain, by Michael, on the software release 94e3f01:  the configuration by
  \h{larpix-lib} through \h{pacman\_cmdserver}, the menu's hello replies seen on felix through the data server. \\
2026-11-04 & 1v5 & Bring-up from felix alone (then \h{agents/bringup.sh}):  tile 1 power (VDDA 2367~mV, VDDD
  1302~mV), full reset, 14 configuration writes, read-back through the data server.  One read at a time exact;  four
  in a row, only the first and last replies:  the data server's PUB socket dropped them (send high-water mark 1). \\
2026-11-04 & 1v5 & \h{PUB\_HWM} 1000:  four reads in a row 4 of 4, 64 in a row 64 of 64;  \h{pacman\_push} 20000
  words, all counted (6680 with HWM 1).  In software 3.6.2. \\
2026-11-05 & 1v5 & The closed loop (\h{AsicRegisterIface}, \h{PacmanAsicRegister} then):  4 reads answered in 6~ms
  (was a fixed 0.5~s subscription sleep);  \h{nrep} 3 without discrepancy;  64 registers in one transmission in
  8~ms;  fire and forget acknowledged with no reply collected and nothing stale after it;  pixel trim 0 written and
  read back, then restored;  \h{init\_network.py} through \h{AsicConfig}, configuration unchanged;  20 of 20 runs. \\
2026-11-05 & 1v5 & Board control (\h{Pacman}):  \h{read\_pacman\_id} 1 after \h{sync\_pacman\_id};  with UART 1
  disabled a read is acknowledged with no reply (\h{complete} False), answered again once enabled;  an internal
  reset of tile 0 keeps every register;  a full reset silences chip 11 and \h{init\_network.py} restores it from
  chip 1;  tile 0 power off and on:  VDDA 2365 $\to$ 29 $\to$ 2367~mV, VDDD 1302 $\to$ 8 $\to$ 1302~mV. \\
2026-11-05 & 1v5 & \h{hello\_asic.py}:  PASS warm (14 of 14 registers, chip ID included);  after \h{reboot}, PASS
  29~s after the reboot (the board answered after about 22~s);  after a supply power cycle (halt, \h{OUTP 0},
  \h{OUTP 1}), PASS 47~s after power-on.  With \h{timeout\_ms} 200 and the swept settles:  0.18--0.19~s, warm and
  from a powered-off tile (6 of 6). \\
2026-11-05 & 1v5 & \h{dump\_registers.py}:  256 registers in 0.10~s, three reads each without a discrepancy;  the
  reset values differing from the datasheet are recorded in \h{larpix.tex}.  More than about 100 reads in one
  transmission lost the rest:  \h{pacman\_cmdserver} moved at most 100 TX buffers per request and the rest waited
  for the next request, its replies arriving late (\h{zipline.tex}, software 3.6.3);  \h{WORDS\_PER\_MESSAGE} 100
  meanwhile, 16000 once the fix was on the bench:  the 256-register scan in one request, 13--15~ms. \\
2026-11-05 & 1v5 & Pedestals (software 3.6.3 candidate):  100 data packets/s from chip 11, all trigger type
  periodic, all 64 channels evenly (the rolling trigger steps through the channels:  about 1.6/s each);  timestamps
  step by 100001 cycles (the period + 1) and fit in 28 bits;  ADC mean about 257 (warm), range 248--542;  the
  occupancy view live.  The data packet layout (datasheet table 16) holds on real packets. \\
2026-11-05 & 1v5 & \h{pacman\_dataserver} quit during the evening (cause not seen:  its output went to
  \h{/dev/null});  restarting it alone brought neither replies nor data back, restarting both servers
  (\h{pacman\_cmdserver} first) did, the ASIC configuration and its running trigger untouched.  Both servers now
  log to \h{/tmp}. \\
\bottomrule
\end{tabularx}
\end{center}

\section{Hydra strands on the 5x5}

Single hydra on the 5x5 (2026-11-07, Michael's scope:  are larpix-lib's configurations sufficient?).  The
UC Davis 5x5 tile on PACMAN tile 1, the 1v5 on software 3.6.3;  \h{demos/hydra\_5x5.py}:  one power-up and full
reset, each strand configured root outward, every register written read back, then 10\,s of pedestals.  The
roots take larpix-lib's own fragments;  every other chip also \h{init\_link\_5x5} (network\_5x5.py's link settings,
on all four ports, since \h{init\_io} has no port).  Hydra files \h{config/hydra/5x5\_io1}--\h{4.yaml} (rows;  row 5
south of 41) and \h{5x5\_one\_strand.yaml} (row 1 east from 11, each column south from it).

\begin{center}\small
\begin{tabularx}{\textwidth}{lL}
\toprule
Step & Result \\
\midrule
5 chips, row 11--15, io channel 1 & configured;  pedestals about 100 packets/s on every chip \\
25 chips, one strand from 11 & configured in 0.3\,s;  420 of 422 reads match;  pedestals about 100/s on all 25 \\
25 chips, four roots, io channels 1--4 & pedestals about 100/s on all 25 (chip 13:  87);  1--2 reads missed per strand \\
\bottomrule
\end{tabularx}
\end{center}

Findings:
\begin{itemize}
\item \h{pedestal\_channels} (channel mask 0) leaves every channel's own trigger on:  on the pixel tile chip 11 sent
  140\,000 packets/s and saturated its UART.  \h{pedestal\_channels\_5x5} masks them (channel mask 1, periodic
  trigger mask 0, global threshold 255), as \h{enable\_pedestal\_5x5.py}.
\item With VDDA off (as \h{pedestal.py} writes the pedestal fragments) the root takes writes, the chips beyond it
  do not:  writes are not passed from chip to chip.  \h{hydra\_5x5.py} keeps VDDA on.
\item Register dumps against network\_5x5.py (chips 11--13):  larpix-lib leaves \h{i\_rx\_clk}, \h{i\_rx\_rst},
  \h{i\_rx\_ext\_trig} at 0 and \h{reset\_threshold} at 255, the chips' power-on values (the spec's reset values:
  2, 2, 2, 0);  larpix-control writes every register (3, 3, 3, 0).  Configuration and pedestals work either way.
\end{itemize}
Open:  in the readback batch, the chip ID read of a root or its first child sometimes has no reply (once
\h{0x04}), and the farthest chip's last read can time out;  the same reads alone always answer.  The data server
quit three times during the session;  restarting both servers brought it back each time.

\section{Student work}

David Yang (\h{dy241/larpix-lib}, 2026-06 to 08), each part as his own commit, authored with his GitHub no-reply
address:  his reply correlator and \h{AsicSpec} read-side fixes (fea7566;  an authorship commit, not in a working
state on its own), adapted to \h{AsicRegisterIface} in the next commit;  his data packet layout and
\h{valid\_data\_packet} (da768fc).  His ASIC simulator with multi-chip grids (\h{sim/}) is not taken yet;  it would
implement the three board interfaces.

\section{Open}

\begin{itemize}
\item The HDF5 data server (tabled):  separate from \h{larpix\_control}, larpix-control's raw format suggested,
  the reference version open.
\item The trigger period as a fragment parameter (needs a number split across a byte array).
\item Why \h{pacman\_dataserver} quit on 2026-11-05 (now logged).
\item Board control:  the power monitors and other replies (deferred:  pokes return nothing for now);  whether
  \h{power\_on} alone leaves the ASICs as after a full reset (untested:  \h{hello\_asic.py} always resets).
\item Deployment variations and more than one channel or board:  the binding of io groups and channels to strands.
\item For 3.7:  the XPUB subscription confirmation, separate config and data streams, polling for reset complete.
\end{itemize}

\end{document}
